\documentclass[12pt]{article}
\usepackage[utf8]{inputenc}
\usepackage[T1]{fontenc}
\usepackage{hyperref}
\title{Wolane Language Resources for K-12 Curriculum Development}
\author{Compiled by Claude Sonnet 5.5 (claude-sonnet-5-5) for Andr\'as Kornai}
\date{}
\begin{document}
\maketitle

\section{Introduction}
Wolane (ISO 639-3 wle) is an East Gurage Ethiosemitic (Afro-Asiatic) language of central Ethiopia, closely related to Silt'e, Zay and Harari. Sources found in search differ on speaker numbers: the grammar's publisher estimates about 70,000 speakers, while Wikipedia gives about 160,000 for the ethnic community. Younger speakers are mostly multilingual and use Amharic as a second language. Searches found no standard writing system and no school materials. Wolane is described as an oral culture.

\section{Primary Educational Resources}
No primers, textbooks or school materials in Wolane were found in search results.

\section{Linguistic and Reference Materials}
\begin{itemize}
\item Ronny Meyer, \emph{Wolane. Descriptive Grammar of an East Gurage Language (Ethiosemitic)}, a review of the grammar: \url{https://journals.sub.uni-hamburg.de/aethiopica/article/view/172}
\item Publisher page for the Meyer grammar (R\"udiger K\"oppe Verlag): \url{https://www.koeppe.de/titel_wolane}
\item Addis Ababa University MA thesis (2011) containing a grammatical sketch of Wolane and documentation of Wolane riddles: \url{https://etd.aau.edu.et/bitstreams/00ebee1d-636b-442d-a40c-1dfb7bb8f4e3/download}
\end{itemize}

\section{Religious and Community Resources}
\begin{itemize}
\item Joshua Project profile of the Wolane people (ethnographic summary only; no scripture or religious materials were found): \url{https://joshuaproject.net/people_groups/19814/ET}
\item Wikipedia, Wolane people: \url{https://en.wikipedia.org/wiki/Wolane_people}
\end{itemize}

\section{Digital and Audio Resources}
No audio, video, corpus or software resources for Wolane were found in search results.

\section{Curriculum Development Context}
Very little exists beyond one descriptive grammar and one thesis. Wolane has no standard orthography, and Amharic is the usual second language. Any curriculum would need orthography development and materials creation with the community, building on Meyer's grammar and the riddle documentation.
\end{document}